\chapter{Evaluation, Limitations, and Evolution}

\section{Introduction}

This chapter evaluates JANUS as a defensive prototype rather than as a
finished commercial product. It separates demonstrated capabilities from
format-dependent behavior and from future work. This distinction is essential
because a deception system loses value when its dashboard reports certainty
that the underlying sensors cannot provide.

\section{Evaluation against the objectives}

\begin{table}[H]
\centering
\caption{Evaluation of the project objectives}
\label{tab:objective-evaluation}
\small
\begin{tabular}{p{0.26\linewidth}p{0.43\linewidth}p{0.21\linewidth}}
\toprule
Objective & Evidence & Status \\
\midrule
Generate varied business decoys & DOCX, XLSX, CSV, JSON, YAML, ENV, and ZIP
artifacts; remote LLM or local fallback & Achieved \\
Prevent unsafe deployment & forbidden-pattern scan, coherence gate,
format inspection, and project-root path policy & Achieved \\
Detect local interaction & inherited SACL, event 4663, Wazuh rules, and
registry correlation & Achieved on Windows \\
Add off-host detection & official Office and AWS token policies with explicit
activation conditions & Partially achieved \\
Preserve evidence & raw alert JSON, Wazuh identifiers, canonical SHA-256, and
normalized signal & Achieved \\
Present honest results & separate local/remote views and explicit confidence
and suppression fields & Achieved \\
Operate as a production SOC platform & high availability, immutable storage,
distributed agents, and production TLS & Out of scope \\
\bottomrule
\end{tabular}
\end{table}

\section{What can be learned about an actor}

The two principal sensors observe different boundaries. Wazuh provides strong
host-side evidence when the file is touched on an enrolled and correctly
audited endpoint. Canarytokens provides a network-side signal after a
format-specific external action. Neither sensor alone identifies a physical
person.

\begin{table}[H]
\centering
\caption{Actor-related information and evidential quality}
\label{tab:actor-evidence}
\small
\begin{tabular}{p{0.24\linewidth}p{0.25\linewidth}p{0.41\linewidth}}
\toprule
Information & Source & Correct interpretation \\
\midrule
Local account and logon identifier & Windows/Wazuh & Direct evidence for the
audited host event; an account may be compromised. \\
Process name and identifier & Windows/Wazuh & Direct evidence for the process
that exercised the audited access right. \\
Object path, access mask, and time & Windows/Wazuh & Direct evidence for the
object and right used at that time. \\
Remote session address or share & 4624/5145 context & Available only when a
compatible session event exists and correlation is justified. \\
Public source IP and time & Canarytokens & Network egress observed by the
provider; a NAT, proxy, or VPN may hide the device. \\
User-Agent and token channel & Canarytokens & Client-claimed protocol metadata;
useful but spoofable. \\
Approximate location and ASN & Provider IP enrichment & Derived from IP
databases, not GPS and not proof of a person's location. \\
MAC address across the Internet & Neither & Unavailable across routed networks
and therefore never generated by JANUS. \\
Remote download from one 4663 event & Neither alone & A read is observed; a
download requires additional SMB, process, or network evidence. \\
Operating system of an unmanaged actor & Neither reliably & A User-Agent may
suggest a client family, but JANUS must label this as an inference. \\
\bottomrule
\end{tabular}
\end{table}

\section{Security and privacy analysis}

JANUS reduces risk through four design choices. First, it never asks the LLM to
generate a real credential and keeps local paths and token secrets out of the
prompt. Second, administrative endpoints require an API key and sensitive
configuration remains outside version control. Third, the generation pipeline
rejects unsafe paths and malformed artifacts before deployment. Fourth, the
evidence layer preserves the provider payload instead of replacing it with an
unverifiable summary.

The same telemetry can contain personal or operational data, including account
names, paths, public IP addresses, and process details. A real deployment must
therefore define access control, retention, lawful purpose, redaction, and
incident-handling rules. Public screenshots must obscure API keys, AWS token
material, email addresses, public IP addresses when unnecessary, and personal
directory names.

\section{Technical limitations}

\subsection{Endpoint visibility}

Wazuh observes only hosts that forward the relevant events. If a HoneyDoc is
copied to a completely unmanaged computer, the original endpoint can show the
access or copy-related context available before departure, but it cannot report
later local activity on the external computer.

\subsection{Office behavior}

Office token activation depends on external-content policy, application
version, Protected View, network controls, and user choices. DOCX opening can
therefore create a reliable local Wazuh event without producing a remote
callback. JANUS does not use macros to bypass this policy.

\subsection{Structured documents}

A JSON, YAML, CSV, or ENV file cannot autonomously make a network request while
remaining passive data. Such a file can still be a useful local decoy. A web
breadcrumb activates only if an actor follows or uses it. AWS credentials
activate only when presented to an AWS API, and the provider may report the
event after a delay \cite{canaryAws}.

\subsection{Laboratory infrastructure}

The validated deployment uses one Windows workstation, a local SQLite database,
and a Docker single-node Wazuh stack. Indexer TLS verification is relaxed for a
local self-signed laboratory certificate. This setting, local process launching,
and SQLite storage must be redesigned before production use.

\section{Threats to validity}

The current experiment demonstrates functional correctness but does not measure
detection probability over a representative attacker population. The observed
Canarytokens IP originated from a controlled test, not from an adversary. File
access performed by PowerShell proves the sensor path, not the behavior of every
Office version. Furthermore, geolocation and User-Agent fields are provider
enrichments whose accuracy is outside JANUS control. Repeated experiments on
separate endpoints and networks are needed before drawing statistical
conclusions.

\section{Proposed evolution}

The next release should preserve the simple, explainable core while improving
deployment quality:

\begin{enumerate}[label=\arabic*.]
  \item expose a public HTTPS callback endpoint with authentication and replay
  protection so remote token events enter the same evidence registry;
  \item enrol a separate Windows victim virtual machine and a Linux endpoint to
  test the trust boundary rather than using only the development host;
  \item enable production certificate verification and secret management;
  \item export JSON/CSV evidence bundles and add role-based dashboard access;
  \item define retention and redaction policies and move evidence to an
  append-only or externally protected store;
  \item measure time-to-detect, false-positive rate, callback rate by format,
  and correlation coverage across repeated scenarios;
  \item add signed deployment manifests and reproducible release packaging.
\end{enumerate}

Experimental fake SSH/HTTP services and prompt-injection or credential-trap
layers should remain outside the core evaluation unless a future research
question explicitly needs them. They are not required for the HoneyDoc plus
Wazuh plus Canarytokens contribution.

\section{Conclusion}

JANUS achieves its academic MVP objective: it creates controlled decoys,
observes a real Windows interaction through Wazuh, supports official remote
token mechanisms, and preserves evidence without inventing attacker
attributes. Its most important future work is not to add more speculative
fields, but to validate the same evidence model on isolated endpoints and to
integrate authenticated remote callbacks.
